\haddockmoduleheading{Data.Ratio}
\label{module:Data.Ratio}
\haddockbeginheader
{\haddockverb\begin{verbatim}
module Data.Ratio (
    Ratio, Rational, (%), numerator, denominator, approxRational
  ) where\end{verbatim}}
\haddockendheader

\begin{haddockdesc}
\item[\begin{tabular}{@{}l}
data Ratio a
\end{tabular}]
{\haddockbegindoc
Rational numbers, with numerator and denominator of some \haddockid{Integral} type.\par
Note that \haddockid{Ratio}'s instances inherit the deficiencies from the type
 parameter's. For example, \haddocktt{Ratio Natural}'s \haddockid{Num} instance has similar
 problems to \haddockid{Natural}'s.\par}
\end{haddockdesc}
\begin{haddockdesc}
\item[\begin{tabular}{@{}l}
instance Integral a => Enum Ratio a\\instance (Storable a, Integral a) => Storable Ratio a\\instance Integral a => Num Ratio a\\instance (Integral a, Read a) => Read Ratio a\\instance Integral a => Fractional Ratio a\\instance Integral a => Real Ratio a\\instance Integral a => RealFrac Ratio a\\instance Show a => Show Ratio a\\instance Eq a => Eq Ratio a\\instance Integral a => Ord Ratio a
\end{tabular}]
\end{haddockdesc}
\begin{haddockdesc}
\item[\begin{tabular}{@{}l}
type Rational = Ratio Integer
\end{tabular}]
{\haddockbegindoc
Arbitrary-precision rational numbers, represented as a ratio of
 two \haddockid{Integer} values.  A rational number may be constructed using
 the \haddockid{{\char '45}} operator.\par}
\end{haddockdesc}
\begin{haddockdesc}
\item[\begin{tabular}{@{}l}
({\char '45}) :: Integral a => a -> a -> Ratio a
\end{tabular}]
{\haddockbegindoc
Forms the ratio of two integral numbers.\par}
\end{haddockdesc}
\begin{haddockdesc}
\item[\begin{tabular}{@{}l}
numerator :: Ratio a -> a
\end{tabular}]
{\haddockbegindoc
Extract the numerator of the ratio in reduced form:
 the numerator and denominator have no common factor and the denominator
 is positive.\par}
\end{haddockdesc}
\begin{haddockdesc}
\item[\begin{tabular}{@{}l}
denominator :: Ratio a -> a
\end{tabular}]
{\haddockbegindoc
Extract the denominator of the ratio in reduced form:
 the numerator and denominator have no common factor and the denominator
 is positive.\par}
\end{haddockdesc}
\begin{haddockdesc}
\item[\begin{tabular}{@{}l}
approxRational :: RealFrac a => a -> a -> Rational
\end{tabular}]
{\haddockbegindoc
\haddockid{approxRational}, applied to two real fractional numbers \haddocktt{x} and \haddocktt{epsilon},
 returns the simplest rational number within \haddocktt{epsilon} of \haddocktt{x}.
 A rational number \haddocktt{y} is said to be \emph{simpler} than another \haddocktt{y'} if\par
\vbox{\begin{itemize}
\item
\haddocktt{\haddockid{abs} (\haddockid{numerator} y) <= \haddockid{abs} (\haddockid{numerator} y')}, and\par
\item
\haddocktt{\haddockid{denominator} y <= \haddockid{denominator} y'}.\par
\end{itemize}}
Any real interval contains a unique simplest rational;
 in particular, note that \haddocktt{0/1} is the simplest rational of all.\par}
\end{haddockdesc}